\documentclass[11pt, conference, onecolumn]{IEEEtran}
\linespread{1.5} % Adds a bit of line spacing to improve single-column readability and increase length

\usepackage{amsmath,amsfonts,amssymb}
\usepackage{booktabs}
\usepackage{tabularx}
\usepackage{geometry}
\usepackage{graphicx}
\usepackage{cite}
\usepackage{hyperref}
\geometry{a4paper, margin=1in}

\begin{document}

\title{Mitigating Sycophancy in LLM Hypothesis Generation: A Cost-Plausibility Analysis of RAG vs. Heterogeneous Multi-Agent Debate}

\author{
    \IEEEauthorblockN{MSE2026005 - Sarvagna Khatri}
    \IEEEauthorblockA{Target: IEEE/ACM Publications\\
    (Double-Blind Review)}
}

\maketitle

\begin{abstract}
While Retrieval-Augmented Generation (RAG) improves the factual grounding of Large Language Models (LLMs), generating plausible scientific hypotheses often requires iterative logical refinement. Multi-Agent Debate (MAD) architectures are theorized to mitigate parametric sycophancy by forcing models into adversarial correction, yet their token-cost efficiency remains poorly quantified in generative tasks. In this study, we perform a controlled evaluation pitting a single LLM with static context (simulated RAG) against an ungrounded heterogeneous two-agent debate loop. Using a frozen 50-sample split of PubMedQA evaluated by an LLM-as-a-judge ($T=0.0$), we measure both scientific plausibility and cumulative token expenditure. Results indicate that the 2-agent debate yields a statistically significant increase in plausibility ($8.58$ vs. $2.14$, $p < 0.001$). However, this qualitative gain requires an $88.5\times$ increase in token consumption ($294,079$ tokens vs. $3,320$ tokens), resulting in a Plausibility Return on Investment (ROI) collapse from $0.644$ to $0.029$. We conclude that while heterogeneous debate drastically improves hypothesis plausibility, its exponential cost renders it highly inefficient for standard queries, providing a critical empirical baseline for developers operating on strict token budgets.
\end{abstract}

\begin{IEEEkeywords}
Large Language Models, Multi-Agent Debate, Retrieval-Augmented Generation, Sycophancy, Cost-Efficiency, Scientific Discovery.
\end{IEEEkeywords}

\section{Introduction}
The generation of novel, plausible scientific hypotheses is an increasingly important frontier for Large Language Models (LLMs). While modern foundational models possess vast parametric knowledge bases spanning the biomedical sciences, they frequently suffer from a critical failure mode: parametric sycophancy. When tasked with open-ended hypothesis generation, standard single-agent LLMs tend to confidently hallucinate mechanisms or conform to early logical fallacies rather than rigorously challenging their initial assumptions \cite{wynn2025debate}.

To address these hallucinations, the community has largely adopted two architectural pathways. The first is Retrieval-Augmented Generation (RAG), which seeks to anchor generative output to retrieved, verifiably true documents \cite{xiong2024kgcoi}. The second is Multi-Agent Debate (MAD), a framework that forces distinct AI personas into iterative adversarial refinement to iteratively correct logical flaws \cite{wynn2025debate}. 

While both architectures have theoretical advantages, there is a critical gap in current literature: the empirical measurement of the cost-to-plausibility Return on Investment (ROI). Specifically, it remains unquantified whether the immense token expenditure of a multi-turn, multi-agent debate is justified when compared to a computationally cheap, static RAG pipeline.

This paper provides an explicit, zero-shot comparative measurement of Cost vs. Plausibility between a static RAG architecture and an ungrounded MAD architecture. By holding the underlying base model constant, we isolate the architectural impact on both scientific plausibility and cumulative token expenditure.

\section{Related Work}

\subsection{Factual Grounding via RAG}
Retrieval-Augmented Generation is currently the industry standard for factual calibration. Recent work by Xiong et al. \cite{xiong2024kgcoi} demonstrates that augmenting LLMs with Knowledge Graphs and retrieved text significantly reduces hallucinations. However, RAG heavily relies on the assumption that the necessary facts already exist in the retrieved literature. For hypothesis generation---a task requiring the synthesis of new ideas---pure retrieval is often insufficient.

\subsection{Multi-Agent Reasoning and Sycophancy}
Multi-Agent Debate (MAD) attempts to simulate human peer review. Wynn et al. \cite{wynn2025debate} explored the failure modes of such systems, noting that while debate can improve reasoning, it frequently suffers from ``conformity'' (where agents sycophantically agree with each other) and extreme token costs. Our study builds on this by explicitly testing a \textit{heterogeneous} persona design (Proposer vs. Challenger) to forcibly break conformity, while rigorously tracking the token cost implications.

\subsection{Automated Evaluation of Hypotheses}
Evaluating the quality of open-ended scientific hypotheses manually is prohibitively slow and subjective. Guo et al. \cite{guo2025ideabench} introduced the IdeaBench framework, demonstrating that carefully prompted LLM-as-a-judge systems can reliably score the plausibility and novelty of generated text. We adapt a constrained plausibility evaluation prompt inspired by this methodology to ensure scalable, reproducible scoring.

\section{Problem Statement and Research Questions}

Our core objective is to determine whether the iterative conversational refinement of an ungrounded two-agent debate loop justifies its exponentially higher token expenditure compared to a single retrieval-augmented LLM.

\textbf{RQ1: Plausibility vs. Architecture.} How does the plausibility of scientific hypotheses generated by a single LLM with static context compare to those generated by an ungrounded two-agent debate system? We hypothesized that the RAG pipeline would produce hypotheses of equal or greater plausibility because domain-specific grounding would outweigh the benefits of adversarial refinement.

\textbf{RQ2: Cost-Efficiency (Token ROI).} What is the token cost required to achieve a baseline plausibility score using single-LLM RAG versus ungrounded debate? We hypothesized that the single-LLM RAG pipeline would achieve its scores using at least 50\% fewer cumulative API tokens than the debate system.

\section{Methodology}

\subsection{Generative Architectures}
We evaluated three distinct generative architectures. For all architectures, we utilized Google's \texttt{gemini-3.5-flash-lite} model to ensure cost constraints and API availability.

\subsubsection{Zero-Shot Single LLM (Baseline)}
A baseline query was issued to a single LLM persona acting as a biomedical expert. No external context or iterative refinement was provided.

\subsubsection{Static Context RAG}
To simulate a deterministic RAG environment, a single LLM was provided with a hardcoded context string containing the exact biomedical abstract associated with the query. This isolated the model's ability to synthesize information from a perfect retrieval hit without introducing the variance of an active embedding database.

\subsubsection{Heterogeneous 2-Agent Debate (Ungrounded)}
We designed a Python \texttt{while}-loop pitting two distinct personas against each other:
\begin{itemize}
    \item \textbf{The Proposer:} Prompted to formulate a highly plausible, scientifically accurate hypothesis answering the query.
    \item \textbf{The Challenger:} Prompted to act as a rigorous peer reviewer, tasked explicitly with identifying logical fallacies, sycophantic assumptions, and lack of evidence.
\end{itemize}
The Proposer generates an initial output. The Challenger critiques it. The Proposer then ingests the critique and revises the hypothesis into a final consensus output. This architecture was \textit{ungrounded}, meaning neither agent had access to external literature or the RAG context string.

\subsection{Resilience Engineering and Token Tracking}
Multi-agent systems executing batch evaluations frequently trigger server rate limits. To prevent catastrophic failure during empirical data collection, we wrapped all API generation calls in an exponential backoff algorithm. Upon encountering a 503 (Server Unavailable) or 429 (Rate Limit) HTTP error, the system automatically paused for $2^{(attempt+1)} + \text{jitter}$ seconds. Furthermore, we utilized the \texttt{usage\_metadata} attribute of the SDK to precisely track every input and output token consumed across all turns of the debate loop.

\section{Experimental Setup}

\subsection{Dataset}
We extracted a strictly frozen random split of 50 samples from the \texttt{qiaojin/PubMedQA} dataset (Seed = 42). The dataset provides medical queries alongside human-verified long-form reference answers, ensuring a high-quality ground truth for the evaluation phase. 

\subsection{LLM-as-a-Judge Evaluation}
To answer RQ1, we employed a separate instance of \texttt{gemini-3.5-flash-lite} as an automated judge. The judge was provided with the original query, the human-verified reference answer, and the generated hypothesis. It was instructed via a standardized system prompt to output a single integer from 1 to 10 rating the scientific plausibility of the generated text. To eliminate stochastic variance during scoring, the evaluation temperature was strictly set to $T=0.0$. 

\section{Results and Discussion}

\subsection{RQ1 Findings: Plausibility}
Table \ref{tab:results} summarizes the architectural performance. The Zero-Shot baseline struggled significantly, averaging a plausibility of $1.54$. The Static Context RAG improved this to $2.14$, demonstrating the expected benefit of factual grounding.

However, the ungrounded Heterogeneous 2-Agent Debate vastly outperformed both pipelines, achieving a mean plausibility of $8.58$. A Wilcoxon Signed-Rank Test between the RAG and Debate scores yielded $W = 0.0$ and a $p$-value of $1.58 \times 10^{-9}$. This conclusively rejects the null hypothesis and contradicts our initial assumption: adversarial refinement (even ungrounded) is mathematically vastly superior to simple single-shot factual grounding for hypothesis generation.

\subsection{RQ2 Findings: The Token Cost Explosion}
While the debate architecture succeeded in maximizing plausibility, it triggered a catastrophic collapse in computational efficiency. Multi-turn architectures inherently scale token usage quadratically per turn, as the entire conversational history must be re-ingested with each API call.

The Zero-Shot and RAG pipelines consumed $2,502$ and $3,320$ tokens respectively across all 50 queries. In stark contrast, the 2-Agent Debate loop consumed a staggering $294,079$ tokens---an $88.5\times$ increase over the RAG baseline. 

\begin{table}[h]
\centering
\caption{Empirical Comparison of Plausibility and Token Cost}
\label{tab:results}
\begin{tabularx}{\columnwidth}{l r r r}
\toprule
\textbf{Architecture} & \textbf{Plausibility} & \textbf{Total Tokens} & \textbf{ROI} \\
\midrule
Zero-Shot LLM & $1.54 \pm 2.18$ & $2,502$ & $0.615$ \\
Static Context RAG & $2.14 \pm 2.95$ & $3,320$ & $0.644$ \\
2-Agent Debate & $\mathbf{8.58} \pm 2.81$ & $\mathbf{294,079}$ & $\mathbf{0.029}$ \\
\bottomrule
\end{tabularx}
\end{table}

We define the Plausibility Return on Investment (ROI) metric to quantify efficiency:
\begin{equation}
\text{ROI} = \frac{\text{Mean Plausibility}}{\text{Total Tokens} / 1000}
\end{equation}

As shown in Table \ref{tab:results}, the ROI for the debate pipeline plummeted to $0.029$. This supports our hypothesis regarding cost inefficiency, confirming that multi-agent debate is not a scalable alternative to standard RAG under strict token or latency budgets.

\section{Limitations and Future Work}

\subsection{Limitations}
\begin{enumerate}
    \item \textbf{Dataset Scope:} The evaluation was constrained to $N=50$ biomedical queries from PubMedQA. While this yielded high statistical significance, the findings may not generalize uniformly to other domains like abstract mathematics or software generation.
    \item \textbf{Evaluation Bias (Self-Preference):} Using the Gemini family to evaluate Gemini-generated text introduces potential intra-model sycophancy. Furthermore, LLM judges often exhibit a ``verbosity bias,'' favoring the inherently longer outputs of the debate loop over the concise outputs of the single-shot RAG.
    \item \textbf{Parametric Leakage:} We labeled the debate architecture ``ungrounded.'' However, because PubMed abstracts are deeply embedded in standard LLM pre-training data, the agents may have been recalling memorized facts rather than demonstrating pure zero-shot logical deduction.
\end{enumerate}

\subsection{Future Work (Ablation Strategy)}
The current experimental data proves that the ungrounded debate outperforms RAG, but it leaves one critical question unresolved: Is the $8.58$ plausibility score driven by the \textit{heterogeneous} adversarial nature of the personas (Proposer vs. Challenger), or is it merely an artifact of allowing a model to generate exponentially more tokens across a loop? 

In future work, we will conduct a controlled ablation study comparing the Heterogeneous debate pipeline against a Homogeneous debate pipeline (two identical Proposer personas). If the Homogeneous pipeline suffers from ``sycophancy collapse'' (yielding lower plausibility scores despite the same token budget), it will definitively prove that adversarial friction is the active ingredient required for high-quality scientific hypothesis generation.

\section{Conclusion}
This research explicitly quantifies the severe trade-off between output quality and computational cost when shifting from retrieval-augmented generation to multi-agent reasoning. We demonstrated that a heterogeneous two-agent debate drastically mitigates parametric sycophancy, increasing the plausibility of scientific hypotheses by over $4\times$ compared to a static RAG baseline. 

However, this qualitative leap comes at an unsustainable computational cost, requiring an $88.5\times$ increase in token expenditure and resulting in a near-zero Plausibility ROI ($0.029$). For production architectures with strict latency and financial constraints, single-LLM RAG remains the only viable heuristic. Unlocking the immense reasoning benefits of Multi-Agent Debate at scale will require urgent algorithmic innovations in dynamic early-stopping and targeted context pruning to control the token-cost explosion.

\bibliographystyle{IEEEtran}
\bibliography{references}

\end{document}
